\section{The remaining three-row generating functions}\label{app:threegf}

This appendix completes the proof of Theorem~\ref{thm:threegf}. It uses the letters, the pair
set $F$, the windows and the diamond test of Section~\ref{sec:three}, and the polynomials
\[
R=1-3z,\qquad L=1-2z-2z^2,\qquad N=1+2z^2,\qquad D=LN,
\]
where $D=1-2z-4z^3-4z^4$ by multiplying out. A \emph{restricted run} is a string of singletons with no two letters $b$ at distance two,
and the \emph{outer} letters are $a$ and $c$.

\subsection{The forbidden windows of the classes containing \texorpdfstring{$B$}{B}}

Lemmas~\ref{lem:dichotomy} and~\ref{lem:sixofnine} and Step~1 of the proof of
Proposition~\ref{prop:threeGAC} list every window that fails in the class $\{A,C\}$. The next
lemma does the same for the ordered class $AB$ and for the triple class. The class $BA$ is the
reversal of $AB$, and the classes $CB$ and $BC$ are the reflections of $AB$ and $BA$ in the
middle row, so Table~\ref{tab:threewindows} covers every class.

\begin{table}[t]
\centering\small
\begin{tabular}{@{}lll@{}}
\toprule
class & windows & forbidden \\
\midrule
$\{A,C\}$ & singletons only & none \\
 & one double & $bAb$, $bAc$, $cAb$, $bCb$, $bCa$, $aCb$ \\
 & two doubles & $bAC$, $cAC$, $ACa$, $ACb$, $aCA$, $bCA$, $CAb$, $CAc$ \\
$AB$ & singletons only & $bab$, $bbb$, $bcb$ \\
 & one double & $bsA$, $Asb$, $bAb$, $bAc$, $cAb$ \\
 & & $aBa$, $aBb$, $bBa$, $cBc$ \\
 & two doubles & $bAB$, $ABa$ \\
$X$ & singletons only & $bab$, $bbb$, $bcb$ \\
 & one double & $bsX$, $Xsb$, $aXa$, $aXb$, $bXa$, $bXc$, $cXb$, $cXc$ \\
\bottomrule
\end{tabular}
\caption{The windows that fail the diamond test, by class, with $s$ an arbitrary singleton.
The class $\{A,C\}$ is from Section~\ref{sec:three}, and the ordered class $AB$ and the class
$X$ are Lemma~\ref{lem:threewindows}. Every window not listed passes.}
\label{tab:threewindows}
\end{table}

\begin{lemma}[forbidden windows in the classes containing $B$]\label{lem:threewindows}
Let $S$ be a candidate.
\begin{enumerate}[(a)]
\item Let $S$ be in the ordered class $AB$, so that $F=\{A,B\}$. Then a window fails if and
only if it is one of the fourteen windows in the rows $AB$ of Table~\ref{tab:threewindows}.
In particular $AsB$ always passes.
\item Let $S$ be in the triple class, so that $F=\{A,B,C\}$. Then a window fails if and only
if it is one of the eleven windows in the rows $X$ of Table~\ref{tab:threewindows}.
\end{enumerate}
\end{lemma}

\begin{proof}
Each window is tested by Lemma~\ref{lem:diamond} on its diamond $D(i,1)$, and three corners
showing means a pass. The singleton windows are Lemma~\ref{lem:dichotomy}(b).

\emph{Step 1. The letter $A$.} A window $sAt$ has the same list as Lemma~\ref{lem:sixofnine}(a),
by the same proof. Only the rival of $bAb$ changes. Its diamond shows $(i+1,3)$ alone, and row
(f) fails it with $x=i+1$, $r=3$ and the row $1\in B_{i+1}$, since $B\in F$. Now let the
window be $Ast$, so the tip $(i,2)$ is black. If $t=b$,
both tips are black. For $s=b$ the diamond shows the vertical pair, which fails by row (c)
since $B\in F$. For $s$ outer it shows only the white point of column $i+1$, and row (f)
fails it, since the rows $1$ and $3$ form $B\in F$. So $Asb$ fails for every $s$. If $t\ne b$,
then for $s=b$ three corners show, and for $s$ outer the diamond shows an oblique pair whose
test in row (e) needs $(i+1,2)$ black, and it is white. So $Ast$ passes for $t\ne b$, and by
reversal $tsA$ fails exactly for $t=b$.

\emph{Step 2. The letter $B$.} In a window $sBt$ column $i+1$ is black. If $s=t=b$ the
presentation is empty, and the window passes by row (b). If $s=b$ and $t\ne b$, the diamond
shows only $(i+2,2)$ with $B_{i+2}=\{t\}$, and row (f) asks whether $\{2,t\}\in F$. So $bBa$
fails and $bBc$ passes, and by the mirror image $aBb$ fails and $cBb$ passes. If $s\ne b$ and
$t\ne b$, the diamond shows the horizontal pair, and row (d) fails it exactly when $s=t$. So
$aBa$ and $cBc$ fail. In a window $Bst$ the tip $(i,2)$ is white, so at most one tip is black.
If $t=b$ and $s=b$, three corners show. If $t=b$ and $s$ is outer, the diamond shows $(i,2)$
and one point of column $i+1$, an oblique pair, and row (e) needs $(i+1,2)$ black, which is
false. If $t\ne b$, three corners show. So $Bst$ and $tsB$ always pass.

\emph{Step 3. The two doubles.} In the window $uAB$ the points $(i+1,1)$ and $(i+1,3)$ are
black and white, the tip $(i+2,2)$ is white, and the tip $(i,2)$ is black exactly when $u=b$.
For $u=b$ the diamond shows the oblique pair $(i+2,2)$, $(i+1,3)$, and row (e) asks whether
$(i+2,3)$ and $(i+1,2)$ are black. Both are, so $bAB$ fails. For $u\ne b$ three corners show.
In the window $ABw$ the tip $(i,2)$ and column $i+1$ are black, and $(i+2,2)$ is black exactly
when $w=b$. For $w=b$ the presentation is empty and the window passes. For $w\ne b$ the
diamond shows only $(i+2,2)$ with $B_{i+2}=\{w\}$, and row (f) asks whether $\{2,w\}\in F$,
which holds exactly for $w=a$. So $ABa$ fails and $ABc$ passes. In the window $AsB$ the tip
$(i,2)$ is black and the tip $(i+2,2)$ is white. For $s=b$ three corners show, and for $s$
outer the diamond shows an oblique pair whose test in row (e) needs $(i+1,2)$ black, and it
is white. So $AsB$ always passes. This proves (a).

\emph{Step 4. The letter $X$.} In a window $sXt$ column $i+1$ is black, so the analysis of
$sBt$ applies with $F=\{A,B,C\}$. The empty case $bXb$ passes, the one-tip case fails for
$t=a$ and $t=c$, since both $\{2,a\}$ and $\{2,c\}$ are in $F$, and the two-tip case fails
exactly when $s=t$. So $bXa$, $bXc$, $aXb$, $cXb$, $aXa$ and $cXc$ fail, and $aXc$ and $cXa$
pass. For $X$ at a tip, the analysis of $Ast$ in Step~1 applies word for word, since it used
only that the tip $(i,2)$ is black and that $B\in F$. So $Xsb$ and $bsX$ fail and every other
such window passes. This proves (b).
\end{proof}

\subsection{Restricted runs and the four-state block}

Read a restricted run letter by letter, and record each letter by a bit, $0$ for an outer
letter and $1$ for $b$. The state after a letter is the pair of bits of the last two letters.
Appending a letter with bit $k$ to the state $(i,j)$ leads to the state $(j,k)$, with weight
$2$ if $k=0$ and weight $1$ if $k=1$. The step is forbidden when $i=k=1$. In the order $00$,
$01$, $10$, $11$ of the states the weighted transition matrix is
\begin{equation}\label{eq:K}
K=\begin{pmatrix}2&1&0&0\\0&0&2&1\\2&0&0&0\\0&0&2&0\end{pmatrix}.
\end{equation}
Let $H=(I-zK)^{-1}$. Its entry $H_{st}$ is the weighted count of the sequences of
transitions from $s$ to $t$ by $z^{\text{length}}$, the empty sequence included. Put
$h=H\mathbf1$ with $\mathbf1$ the all-ones column. So $h_s$ counts all continuations of a run
from the state $s$, the empty one included.

We need four scalar functions of the runs that touch a double, read outward from the double.
For $n\ge1$ let $g_0$ count the restricted runs $w_1w_2\cdots w_n$ in which $w_1$ is a
prescribed outer letter, by $z^n$, and let $g_1$ count those with $w_1=b$. Let $u$ and $v$
count the same runs with the extra condition that $w_2$, if present, is outer. The prescribed
first letter has one choice in $g_0$ and $u$, and every later outer letter has two.

\begin{proposition}[the four-state block]\label{prop:threeblock}
Let $K$ be the matrix~\eqref{eq:K}.
\begin{enumerate}[(a)]
\item The determinant is $\det(I-zK)=D$.
\item In the order $00$, $01$, $10$, $11$ of the states,
\[
h^{\top}=D^{-1}\bigl(1+z+3z^2+2z^3,\,1+z,\,1+2z^2+2z^3,\,1\bigr).
\]
\item The run functions are $g_0=zh_{00}$, $g_1=zh_{01}$, $u=zh_{10}$ and $v=zh_{11}$, that is,
\begin{equation}\label{eq:runs}
\begin{aligned}
g_0&=\frac{z(1+z+3z^2+2z^3)}D,& g_1&=\frac{z(1+z)}D,\\
u&=\frac{z(1+2z^2+2z^3)}D,& v&=\frac zD.
\end{aligned}
\end{equation}
\end{enumerate}
\end{proposition}

\begin{proof}
(a) Expand $\det(I-zK)$ along its first row. The minor of the entry $1-2z$ is $1$, and the
minor of the entry $-z$ is
$\det\begin{pmatrix}0&-2z&-z\\-2z&1&0\\0&-2z&1\end{pmatrix}=-4z^2-4z^3$. So
$\det(I-zK)=(1-2z)+z(-4z^2-4z^3)=D$.

(b) The vector $h$ is the unique solution of $h=\mathbf1+zKh$, that is,
$h_{00}=1+z(2h_{00}+h_{01})$, $h_{01}=1+z(2h_{10}+h_{11})$, $h_{10}=1+2zh_{00}$ and
$h_{11}=1+2zh_{10}$. The last two give $h_{11}=1+2z+4z^2h_{00}$, the second then gives
$h_{01}=1+3z+2z^2+(4z^2+4z^3)h_{00}$, and the first becomes
$h_{00}(1-2z-4z^3-4z^4)=1+z+3z^2+2z^3$. Substituting back gives the other three entries.

(c) A run counted by $g_0$ is its prescribed first letter, with weight $z$, followed by an
arbitrary continuation. A fictitious outer letter before the first letter puts the run in the
state $00$ and restricts nothing, since the letter two places after it is $w_2$, which is
free. So $g_0=zh_{00}$, and in the same way $g_1=zh_{01}$. A run counted by $u$ is its first
letter alone, or its first letter, an outer second letter with two choices, and a
continuation from the state $00$. So $u=z+2z^2h_{00}=zh_{10}$ by the third equation of (b),
and in the same way $v=z+2z^2h_{10}=zh_{11}$ by the fourth. Now (b) gives~\eqref{eq:runs}.
\end{proof}

\subsection{The triple class}

\begin{proposition}[the triple class]\label{prop:GX}
With $u$ and $v$ from~\eqref{eq:runs},
\begin{equation}\label{eq:GX}
G_X(z)=z\bigl(1+4u+2v+2u^2+v^2\bigr).
\end{equation}
Its reduced denominator divides $D^2$.
\end{proposition}

\begin{proof}
A word of the triple class is $UXW$ with $U$ and $W$ strings of singletons, and the letter $X$
contributes $z$. By Lemma~\ref{lem:threewindows}(b), $U$ and $W$ are restricted runs, read
outward from the $X$. The windows $bsX$ and $Xsb$ say that the second letter of each run,
counted from the $X$, is outer if it exists. So a nonempty $U$ is counted by $u$ if its letter
next to $X$ is outer and by $v$ if it is $b$. The same holds for $W$. If both runs are
empty the word is $X$, which contributes $1$ inside the parentheses. If exactly one run is
nonempty, it contributes $2u+v$, and there are two sides. If both are nonempty, the window at
the $X$ allows exactly the pairs $(a,c)$, $(b,b)$ and $(c,a)$ of letters next to it, which
contribute $2u^2+v^2$. Every window of the word is inside a run or is one of the windows
treated, and $u$ and $v$ have the denominator $D$.
\end{proof}

The word $bXb$ of Section~\ref{sec:three} is the case $v\cdot v$ at width $3$. One factor $N$
cancels in $G_X$, so its reduced denominator is $L^2N$, but only the bound $D^2$ is used.

\subsection{The class \texorpdfstring{$\{A,B\}$}{\{A,B\}}}

A word of the ordered class $AB$ is $UAVBW$ with $U$, $V$, $W$ strings of singletons. By
Lemma~\ref{lem:threewindows}(a) the three strings are restricted runs. The windows that cross
a double are $bsA$, $Asb$, $bAb$, $bAc$, $cAb$, $aBa$, $aBb$, $bBa$, $cBc$, $bAB$ and $ABa$.
The window $bsA$ says that $U$, read outward from $A$, has an outer second letter. So a
nonempty $U$ is counted by $u$ or $v$ according to whether its letter next to $A$ is outer or
$b$. The run $W$, read outward from $B$, has no second-letter condition, since $B$ at a tip
never fails. So a nonempty $W$ is counted by $g_0$ or $g_1$ according to its letter next to $B$.

When a singleton $t$ stands directly right of $A$, the window $uAt$ restricts the letter $u$
of $U$ next to $A$. Let $\ell_t$ be the weighted count of the allowed $U$, the empty $U$
included. When $t$ stands directly left of $B$, let $r_t$ be the weighted count of the allowed
$W$, the empty $W$ included. The forbidden windows $bAb$, $bAc$, $cAb$ and $aBa$, $aBb$,
$bBa$, $cBc$ give
\begin{equation}\label{eq:lr}
\begin{aligned}
(\ell_a,\ell_b,\ell_c)&=(1+2u+v,\ 1+u,\ 1+2u),\\
(r_a,r_b,r_c)&=(1+g_0,\ 1+g_0+g_1,\ 1+g_0+g_1).
\end{aligned}
\end{equation}
Indeed next to $At$ every $u$ is allowed for $t=a$, only $u=a$ for $t=b$, and $u\ne b$ for
$t=c$. Next to $tB$ only $w=c$ is allowed for $t=a$, $w\ne a$ for $t=b$, and $w\ne c$ for
$t=c$.

We also need the block $K$ with its final letter marked. Take a start state $s\in\{00,10\}$
and a bit $\beta$. Let $J_{s\beta}$ be the weighted count of the sequences of transitions of
$K$ from $s$ that end in a state whose second bit is $\beta$, the empty sequence included. So
$J$ is formed by the rows $00$ and $10$ of $HE$, where $E$ is the $4\times2$ matrix with
$E_{(i,j),\beta}=1$ exactly when $j=\beta$. The columns $HE_{\cdot\beta}$ solve
$x=E_{\cdot\beta}+zKx$. Solving these two systems as in the proof of
Proposition~\ref{prop:threeblock}(b) gives
\begin{equation}\label{eq:J}
J=\frac1D\begin{pmatrix}1+2z^2+2z^3&z+z^2\\1&2z^2+2z^3\end{pmatrix}.
\end{equation}
The constant term $J(0)=\bigl(\begin{smallmatrix}1&0\\1&0\end{smallmatrix}\bigr)$ records the
empty sequence, and the two rows of $J$ add up to $h_{00}$ and $h_{10}$, as they must.

\begin{proposition}[the class $\{A,B\}$, polynomial form computer-assisted]\label{prop:GAB}
With~\eqref{eq:runs}, \eqref{eq:lr} and~\eqref{eq:J},
\begin{equation}\label{eq:GAB}
\begin{aligned}
G_{AB}(z)={}&z^2(1+2u)(1+g_0+g_1)+z^3(\ell_ar_a+\ell_br_b+\ell_cr_c)\\
&+z^4\,(\ell_a+\ell_c,\ \ell_b)\,J\begin{pmatrix}r_a+r_c\\2r_b\end{pmatrix}.
\end{aligned}
\end{equation}
Moreover,
\begin{equation}\label{eq:PB}
G_{AB}(z)=\frac{P_B(z)}{D^3},\qquad P_B(z)=z^2(1+z)(1+z^2-3z^3+4z^6+4z^7),
\end{equation}
and the remainders of $P_B$ on division by $L$ and by $N$ are
\begin{equation}\label{eq:rem}
P_B\bmod L=\tfrac z2,\qquad P_B\bmod N=\tfrac14-\tfrac z2.
\end{equation}
\end{proposition}

\begin{proof}
We split by the length of the middle run $V$.

\emph{Step 1. The doubles are adjacent.} The word is $UABW$, and the crossing windows are $uAB$
and $ABw$. By Lemma~\ref{lem:threewindows}(a) the letter $u$ of $U$ next to $A$ is not $b$,
and the letter $w$ of $W$ next to $B$ is not $a$. The second-letter condition on $U$ stays.
So $U$ contributes $1+2u$, $W$ contributes $1+g_0+g_1$, and the doubles contribute $z^2$.

\emph{Step 2. One singleton between the doubles.} The word is $UAtBW$, and the window $AtB$
passes. The windows $uAt$ and $tBw$ are counted by $\ell_t$ and $r_t$, and no other window
crosses a double. So the case contributes $z^3\sum_t\ell_tr_t$.

\emph{Step 3. At least two singletons between the doubles.} The word is $UAVBW$ with
$V=v_1v_2\cdots v_k$ and $k\ge2$. The window $Av_1v_2$ is $Asb$ when $v_2=b$, so $v_2$ is
outer. The window $uAv_1$ is counted by $\ell_{v_1}$, the window $v_kBw$ by $r_{v_k}$, and
the window $v_{k-1}v_kB$ passes by Step~2 of the proof of Lemma~\ref{lem:threewindows}.
Inside $V$ the only condition is that $V$ is a restricted run. So the case counts the
restricted runs $v_1v_2\cdots v_k$ with $v_2$ outer, each weighted by
$z^{k+2}\ell_{v_1}r_{v_k}$.

Fix $v_1$ and $v_2$. The run is in the state $(i,0)$ after $v_2$, where $i$ is the bit of
$v_1$, and the letters $v_3,\dots,v_k$ are a sequence of $k-2$ transitions of $K$ from that
state. In $K$ every outer letter carries the weight $2$ for its two colors. If $k=2$, the
last letter is $v_2$ itself, and the sum of $r_{v_2}$ over its two colors is $r_a+r_c$. Let
$k\ge3$. For each sequence of bits the colors of the outer letters are independent. So if
$v_k$ is outer, summing $r_{v_k}$ over its two colors replaces its weight $2$ by $r_a+r_c$,
and if $v_k=b$ the weight is $r_b$, while $v_2$ keeps its two colors. So the weighted count
of all $V$ with a given $v_1$ is
\[
J_{s(v_1),0}\,(r_a+r_c)+J_{s(v_1),1}\,2r_b,
\]
where $s(a)=s(c)=00$ and $s(b)=10$, and the case $k=2$ is the constant term $J(0)$
of~\eqref{eq:J}. Weighting $v_1$ by $\ell_{v_1}$, summing over its three values, and adding
the factor $z^4$ for $A$, $v_1$, $v_2$ and $B$ gives the third term of~\eqref{eq:GAB}. The
three steps are disjoint and exhaustive, and every window of the word is inside a run or is
one of the windows treated, so~\eqref{eq:GAB} is proved.

\emph{Step 4. The polynomial form.} Substituting~\eqref{eq:runs}, \eqref{eq:lr}
and~\eqref{eq:J} into~\eqref{eq:GAB} and clearing denominators is exact polynomial algebra,
and it is computer-assisted. The script \texttt{three\_row\_closed\_form.py} does it with
SymPy and finds~\eqref{eq:PB}. It also checks the first coefficients $7$, $31$, $120$, $449$,
$1622$ of~\eqref{eq:GAB} against the class counts of the data file named in
Section~\ref{sec:three}. It checks the three terms of~\eqref{eq:GAB} separately against those
counts split by the length of $V$. The remainders~\eqref{eq:rem} then follow by hand from
the quotients
\begin{align*}
P_B&=L\bigl(-2z^8-2z^7-z^6+z^4-\tfrac z2\bigr)+\tfrac z2,\\
P_B&=N\bigl(2z^8+4z^7+z^6-2z^5-2z^4+\tfrac32z^2+\tfrac z2-\tfrac14\bigr)\\
&\qquad+\tfrac14-\tfrac z2,
\end{align*}
which are two multiplications.
\end{proof}

\subsection{The total series and its denominator}

\begin{theorem}[the exact denominator, computer-assisted]\label{thm:denominator}
Let $A$ be the series $\sum_{m\ge3}a_3(m)z^m$ of Theorem~\ref{thm:threegf}(a). Then $A$ is a
rational function whose reduced denominator is exactly $Q=R^3L^3N^3$, and each of the five
roots of $Q$ is a pole of order three.
\end{theorem}

\begin{proof}
\emph{Step 1. Divisibility.} By Proposition~\ref{prop:threeGAC}, $G_{AC}$ has the denominator
$R^3$. By Propositions~\ref{prop:GAB} and~\ref{prop:GX}, $G_{AB}$ and $G_X$ have denominators
dividing $D^3$ and $D^2$. So $2G_{AC}+4G_{AB}+G_X-z-12z^2$ has a denominator dividing
$R^3D^3=Q$. The polynomials $R$, $L$ and $N$ are pairwise coprime. Indeed $L-N=-2z-4z^2$ vanishes only at $z=0$ and $z=-\frac12$, where $N\ne0$, and
$R(1/3)=0$ while $L(1/3)=\frac19$ and $N(1/3)=\frac{11}9$.

\emph{Step 2. The pole at $z=1/3$.} The numerator of $G_{AC}$ in~\eqref{eq:GAC} takes the
value $\frac19\bigl(1-\frac43+\frac{10}9-\frac8{27}-\frac1{27}\bigr)=\frac4{81}$ at
$z=\frac13$, so $2G_{AC}$ has a pole of order exactly three there. The functions $G_{AB}$ and
$G_X$ are analytic at $z=\frac13$, since $D(\frac13)\ne0$.

\emph{Step 3. The poles at the roots of $L$ and $N$.} Let $\zeta$ be a root of $L$. Then
$N(\zeta)\ne0$ by Step~1, and by~\eqref{eq:rem} $P_B(\zeta)=\zeta/2\ne0$, since $L(0)=1$. So
$G_{AB}=P_B/D^3$ has a pole of order exactly three at $\zeta$, while $G_X$ has a pole of order
at most two there and $G_{AC}$ is analytic. Let $\zeta$ be a root of $N$, so
$\zeta^2=-\frac12$. Then $L(\zeta)\ne0$ and $P_B(\zeta)=\frac14-\frac\zeta2\ne0$, since
$\zeta\ne\frac12$, and the same argument applies.

So the five roots of $Q$, namely $\frac13$, the two roots of $L$ and the two roots of $N$, are
poles of order three of the series. There are no other poles, and the reduced denominator is
$Q$.
\end{proof}

Steps 1 and 2 are proved by hand. Step~3 rests on the polynomial form~\eqref{eq:PB}, which
is computer-assisted, and on the remainders~\eqref{eq:rem}, which follow from it by hand. The
three cubes in $Q$ have a reading. Each word has three singleton runs, free runs for the class
$\{A,C\}$, which gives $R^3$, and restricted runs for the classes containing $B$, which gives
$D^3$. The triple class has only two runs, so it contributes no new pole.
